% Tension à vide d'une pile Daniell au cours de sa décharge (corrigé du défi D08)
\begin{tikzpicture}
\begin{axis}[otMesures,otAxesBout,otLecture,width=104mm,height=50mm,xmin=0,xmax=1.06,ymin=0,ymax=1.3,
  xtick={0,0.2,...,1},ytick={0,0.2,...,1.2},
  xticklabel style={/pgf/number format/fixed,/pgf/number format/precision=1},
  yticklabel style={/pgf/number format/fixed,/pgf/number format/fixed zerofill,/pgf/number format/precision=1},
  xlabel={fraction $q/Q_{\max}$ de la capacité utilisée},ylabel={tension à vide $U$ (V)}]
\addplot[otNavy,line width=1pt,domain=0:0.9999,samples=400] {1.10-0.0296*log10((1+x)/(1-x))};
\addplot[only marks,mark=*,mark size=1.6pt,otOrange] coordinates {(0.5,1.086)(0.9,1.062)(0.99,1.032)(0.9999,0.973)};
\node[pcLegende,anchor=south west,align=left] at (axis cs:0.04,0.3) {$U$ ne s'annulerait que pour $[\ce{Cu^2+}]\approx\qty{1e-38}{mol.L^{-1}}$,\\c'est-à-dire bien après l'épuisement des ions \ce{Cu^2+}.};
\end{axis}
\end{tikzpicture}
\par\smallskip\legende{La tension à vide reste voisine de 1,1~V presque jusqu'à la fin : la durée de vie est fixée par le réactif limitant.}
